\section{UAV Radar System Feasibility}
\todo{Develop these to more completely engage with what they're doing. Read the works fully}

Sysak et al. \cite{sysak2025_lightweight_compact_pulse_radar_uav_collision_avoidance} present a lightweight X-band monopulse radar designed explicitly for mid-air collision avoidance on medium UAV platforms, addressing the range shortfall that often limits compact frequency modulated continuous wave (FMCW)-based solutions. Their system integrates a multi-channel, forward-looking antenna with an FPGA-based pulse compression on a single 60 gram PCB, enabling online range processing and adaptive CA-CFAR detection, while keeping average power consumption below 80 Watts. Field trials in urban environments demonstrate reliable detection of large obstacles at 1–3 km, providing manoeuvring standoff times on the order of tens of seconds. This work is strongly aligned with UAV Detect-and-Avoid motivations and demonstrates that high-performance sensing can be achieved under strict SWaP constraints with a pulsed, multi-channel architecture. 

However, the evaluation is primarily conducted against large, high radar cross-section infrastructure targets, and does not characterise sensitivity to small airborne targets in realistic encounter geometries. In addition, while the work acts as a proof of concept for UAV radar systems, it does not address multi-target tracking nor the latency, throughput and memory usage of a complete onboard pipeline. These omissions are important for DAA, where hardware is limited, and efficient data handling is as critical as raw detection range. The proposed project complements Sysak et al. by benchmarking detection and tracking sensitivity on small moving targets, and by investigating GPU-oriented implementations and pipeline architectures suitable for low-compute embedded platforms, including the performance impact of memory movement.


\subsection{End-to-end UAV DAA onboard processing pipelines}
Feasibility studies with more emphasis on the radar DAA processing pipelines have been documented. Three that are especially notable include the systems designed by Shi et al. \cite{6712628}, Moses et al. \cite{6044363}, and Newmeyer et al. \cite{7577322}. Shi et al, similarly motivated by the need to detect-and-avoid non-cooperative targets, describe the design of a small FMCW phased array built mainly from off-the-shelf development boards. The design is only suitable for large UAVs proposed project will target smaller platforms. Additionally, the full radar processing was not entirely onboard and online, splitting steps between an onboard FPGA and offboard post processing. On the other hand, the proposed research will design the software for fully onboard, online tracking. 

Moses et al. illustrate the capabilities of a lightweight, single channel FMCW radar. Testing was performed on a quadcopter. A major component of their analysis emphasises identification via micro-Doppler modulation, which is largely evaluated through Doppler-domain signatures over end-to-end range--Doppler detection and tracking. While this is well motivated for classification, it doesn't address the tracking targets in the  range--Doppler space. It also remains unclear how the small aperture and SWaP limits affect range-domain SNR, achievable range resolution, and the resulting false alarm behaviour. The proposed research fills this gap by benchmarking range--Doppler processing, CFAR detection, monopulse angle estimation, and multi-target tracking in low-SNR regimes, and by quantifying how these choices trade sensitivity against latency and throughput on embedded GPU hardware. 

Finally, Newmeyer et al. demonstrated a 120 gram, 8 Watt FMCW phased array radar and processing pipeline that utilised the FPGA fabric on a Zynq ARM CPU. Power, compute cost, latency distributions, and hardware resources were closely monitored, and the project clearly demonstrates that radar DAA systems can operate under strict SWaP constraints. However, the study only evaluates one hardware configuration, and one radar algorithm sequence. The proposed research fills this gap by systematically comparing alternative algorithmic pathways and low-level implementations, and by evaluating how pipeline bottlenecks and end-to-end performance change across different embedded compute and memory configurations, including GPU-based architectures. 

While these works demonstrate early feasibility of small-platform DAA radar, their hardware stacks are not fully representative of modern embedded radar processing platforms in the sense that they do not utilise GPUs or occur entirely onboard, and their evaluations do not explore how algorithmic selection relates to sensitivity to faint target signals. Consequently, there remains a gap in the literature regarding how to integrate and benchmark alternative radar processing and tracking algorithms for small conformal apertures under SWaP constraints, particularly when optimising jointly for sensitivity, latency, and throughput. 
